% Generated draft: requires mathematical and source review.
\documentclass[11pt]{article}
\usepackage[margin=25mm]{geometry}
\usepackage{fontspec}
\setmainfont{Latin Modern Roman}
\usepackage{amsmath,amssymb,mathtools,unicode-math}
\setmathfont{Latin Modern Math}
\usepackage{enumitem,needspace,xurl,hyperref}
\setlength{\parindent}{0pt}
\setlength{\parskip}{5pt}
\setlength{\emergencystretch}{3em}
\newcommand{\N}{\mathbb{N}}
\newcommand{\Z}{\mathbb{Z}}
\newcommand{\Q}{\mathbb{Q}}
\newcommand{\R}{\mathbb{R}}
\newcommand{\C}{\mathbb{C}}
\newcommand{\F}{\mathbb{F}}
\newcommand{\A}{\mathbb{A}}
\newcommand{\PP}{\mathbb{P}}
\newcommand{\E}{\mathbb{E}}
\newcommand{\eps}{\varepsilon}
\newcommand{\dd}{\,\mathrm{d}}
\begin{document}
\section*{UnsolvedMath ID: 800063}
\textbf{Source problem code:} catalogue entry
\subsection{Definitions and mathematical statement}
The supplied question uses the usual terminology of three-dimensional contact topology. To describe that interpretation, let \(M\) be a smooth oriented \(3\)-manifold. The excerpt does not explicitly specify the dimension or orientation of \(M\), whether it is connected, compact or closed, or any boundary conditions; these hypotheses require confirmation from the source.

A cooriented contact structure on \(M\) is a plane field \(\xi=\ker\alpha\), where \(\alpha\) is a smooth nowhere-zero \(1\)-form with \(\alpha\wedge\dd\alpha\) nowhere zero, and the coorientation is represented by positive multiples of \(\alpha\). Relative to the orientation of \(M\), the structure is positive or negative according to the sign of \(\alpha\wedge\dd\alpha\). A contact structure is tight if it contains no overtwisted disk, in the standard contact-topological sense.

Suppose \(\xi_-\) and \(\xi_+\) are cooriented tight contact structures on \(M\), represented by forms \(\alpha_-\) and \(\alpha_+\) satisfying
\[
\alpha_-\wedge\dd\alpha_-<0,
\qquad
\alpha_+\wedge\dd\alpha_+>0.
\]
Suppose also that there exists a vector field \(Z\) on \(M\) such that
\[
Z(p)\notin\xi_-(p)\cup\xi_+(p)
\qquad\text{for every }p\in M,
\]
or equivalently \(\alpha_-(Z)\ne0\) and \(\alpha_+(Z)\ne0\) everywhere. The supplied transversality assumption does not require \(Z\) to be positive with respect to either prescribed coorientation.

Does \(M\) admit a cooriented codimension-one Reebless \(C^0\)-foliation \(\mathcal F\)? Here coorientation means a consistent continuous choice of transverse orientation, and Reebless means that \(\mathcal F\) has no Reeb component, namely no embedded foliated solid torus with the standard Reeb foliation and its boundary torus as a leaf. The precise regularity intended by \(C^0\) is not specified in the excerpt: a topological foliation described by continuous local product charts and a foliation with differentiable leaves and a continuous tangent plane field are different possible conventions. This regularity qualification, together with the missing hypotheses on \(M\), remains unresolved.
\subsection{Short English statement}
If a manifold \(M\) carries tight cooriented contact structures of opposite signs and a vector field transverse to both, must it admit a cooriented Reebless \(C^0\)-foliation? The supplied excerpt leaves the hypotheses on \(M\) and the precise meaning of \(C^0\) unresolved.
\subsection{Sources}
\begin{itemize}
\item Local UnsolvedMath registry, permanent ID \(800063\), source problem code \texttt{catalogue entry}, titled \emph{Foliations and Diffeomorphism Groups --- Question 4}. The supplied record attributes the problem to Nariman and Yazdi, with proposed year \(2025\), and labels it open; that label is supplied registry information, not an independently verified mathematical assessment.
\item Nariman and Yazdi, \emph{Problem List on Foliations and Diffeomorphism Groups} (2025), Question 4. Supplied source URL: \url{https://arxiv.org/abs/2503.18274}. The provenance records access on \texttt{2026-10-16} and extraction from arXiv LaTeX. These are supplied provenance statements; no source visit or independent verification is asserted here.
\end{itemize}
\end{document}
